\documentclass[10pt,a4paper]{amsart}
\usepackage[margin=19mm]{geometry}
\usepackage{amsmath,amssymb,mathtools}
\usepackage[scr=rsfs,cal=euler]{mathalfa}
\usepackage{xfrac}
\usepackage[unicode,hidelinks,bookmarks=false]{hyperref}
\newcommand{\ie}{i.e.\ }
\newcommand{\cf}{cf.\ }
\newcommand{\resp}{respectively\ }
\newcommand{\Subsection}{\subsection}
\providecommand{\nobreakdash}{\nobreak\hbox{-}}
\setlength{\emergencystretch}{2em}
\multlinegap0pt
\usepackage{bm}
\usepackage{bbm}
\usepackage{dsfont}
\usepackage{xspace}
\usepackage{cancel}
\usepackage{mathtools}
\usepackage[shortlabels]{enumitem}
\usepackage{amsthm}
\makeatletter
\@ifundefined{theorem}{\newtheorem{theorem}{Theorem}}{}
\@ifundefined{definition}{\newtheorem{definition}[theorem]{Definition}}{}
\makeatother
\makeatletter
\newcommand{\arena}[1]{\mathbb{#1}}                     % an arena \mathbb{A}
\newcommand{\base}[1]{\mathcal{#1}}                    % a base \mathfrak{B}
\expandafter\ifx\csname bprooftree\endcsname\relax
\newenvironment{bprooftree}
  {\leavevmode\hbox\bgroup}
  {\DisplayProof\egroup}
\fi
\newcommand{\brule}[1]{\triangleright{#1}}         % a premiss pair (P/q)
\newcommand{\deriv}[1]{\vdash_{#1}}       % Sandqvist derivability \supp{\mathfrak{B}}
\newcommand{\fillBox}{\hfill$\Box$}
\newcommand{\toarena}[1]{\mathrel{\textsf{Arena}}(#1)}           % base to arena
\newcommand{\wstrat}[1]{\blacktriangleright_{#1}}            % winning strategy \triangleright_{A}
\makeatother
\title[Inferentialist Game Semantics (Extended Abstract)]{Inferentialist Game Semantics (Extended Abstract)}
\author{Joaquim T. Waddington, Alexander V. Gheorghiu, David J. Pym}
\date{Source: arXiv:2607.13855v2 (CC BY 4.0)}
\begin{document}
\maketitle

\noindent\emph{Source and licence.} Inferentialist Game Semantics (Extended Abstract). Version arXiv:2607.13855v2 (CC BY 4.0), distributed under the Creative Commons Attribution 4.0 International licence, as recorded in the arXiv licence metadata for this exact version and shown on the abstract page https://arxiv.org/abs/2607.13855v2. Full rights notice: licenses/arxiv-2607.13855-CC-BY-4.0.txt.\par\smallskip

\noindent\textit{Editorial scope.} Counted unit: the Theorem whose source label is \texttt{thm:game-soundness} in the article (source0.tex lines 766--771, source label \texttt{thm:game-soundness}), delivered together with the article's own complete original proof of it (source0.tex lines 772--893, one \texttt{proof} environment of 122 lines). No mathematics is written, repaired or re-derived: statement and proof are the article's own text line for line. Required context retained uncounted: def:derivability-in-a-base (lines 271--287); def:play (lines 397--420); def:strategy (lines 424--445). Every definition, display and label of those lines is retained unchanged. Omitted from this package and not redistributed: 1--270, 288--396, 421--423, 446--765, 894--1561. Nothing inside a retained excerpt is omitted or altered: every retained body is delivered exactly as the article's lines are written, so no omission receipt of any kind accompanies this package. Citations: the retained excerpts carry no citation command at all, so no bibliography entry of the article is transcribed and no reference is redistributed. Reference adaptation: none was needed and none was made; every reference inside the delivered unit resolves inside it, so no unresolved reference or citation is produced. Printed slips kept literal: no printed word, symbol, hypothesis or number of the counted statement or of its proof was altered. Layout and class: the article's own class is \texttt{lipics-v2021} with packages including fontenc, amsmath, amssymb, amscd, amsfonts, amsthm, latexsym, esvect, bussproofs, natded, proof, prooftrees, forest, graphicx; the wrapper is the lane's pinned \texttt{amsart} editorial wrapper at 10pt on A4, which restates the theorem-like environments the retained text needs and adds the article's own macro definitions that the retained text needs (\texttt{\textbackslash arena}, \texttt{\textbackslash base}, \texttt{\textbackslash brule}, \texttt{\textbackslash deriv}, \texttt{\textbackslash fillBox}, \texttt{\textbackslash toarena}, \texttt{\textbackslash wstrat}), with extra wrapper packages (bm, bbm, dsfont, xspace, cancel, mathtools, enumitem). The delivered numbering is the unit's own. Assets: no retained span contains a figure environment, a graphics-inclusion command, a PDF-page inclusion, a table or a TikZ picture; no image file is copied into this package, the package declares no asset dependency and no image file is redistributed with it. Licence: version arXiv:2607.13855v2 of this article is distributed under the Creative Commons Attribution 4.0 International licence, as recorded in the arXiv licence metadata for this exact version and shown on the abstract page https://arxiv.org/abs/2607.13855v2. A full rights notice is delivered at licenses/arxiv-2607.13855-CC-BY-4.0.txt. Characters: every retained excerpt is pure ASCII, so the delivered engine, XeLaTeX with its default Latin Modern text font, reports no missing character.
\begin{definition}[Derivability in a base]
\label{def:derivability-in-a-base}
Let $\base{B}$ be a base. The \emph{derivability relation} $\deriv{\base{B}}$,
relating sets of atoms to atoms, is defined
inductively by the following clauses:
\[
\begin{array}{rl}
    \mbox{\rm (Ref)} & \mbox{if $p \in S$, then $S \deriv{\base{B}} p$} \\
    \mbox{\rm (App)} & \mbox{if $(Q_1 \brule q_1), \ldots, (Q_n \brule q_n)
        \brule r$ is a rule of $\base{B}$ and
        $S \deriv{\base{B} \cup Q_i} q_i$, for each $i = 1,...,n$,} \\
        & \mbox{then $S \deriv{\base{B}} r$}
\end{array}
\]
For this to make sense, we identify an atom $p$ with the
$0$th-level rule $\langle p \rangle$. \fillBox
\end{definition}

\begin{definition}[Play]
\label{def:play}
A \emph{play} for an arena $\arena{A}$ is a finite sequence of moves $m_0, m_1, \ldots, m_n$ satisfying the following conditions.
    \begin{enumerate}[label=(\roman*)]
        \item The move $m_0$ is an O-question, called the \emph{initial question}.
        \item There is an index $I \ge 0$ such that $m_1, \ldots, m_I$ are all O-answers,
      called the \emph{initial assumptions} (if $I = 0$ there are none).
        \item After the initial question and the initial assumptions, moves alternate between Opponent and Proponent.
        \item Every root question --- that is, every question whose node is a root of the
      forest of $\arena{A}$ --- is a P-question, and its \emph{justifying question}
      is the initial question.
\item For each question $m_j$ with $j > I$ that is not a root question, there is a
      question $m_k$ with $k < j$ whose node is the immediate predecessor, in
      $\arena{A}$, of the node of $m_j$; we call $m_k$ the \emph{justifying
      question} for $m_j$.
        \item For each answer $m_i$ with $i > I$, there is a question $m_k$ with $k < i$ at the same node of $\arena{A}$. If $m_j$ is the justifying question for $m_k$, we call $m_j$ the justifying question for $m_i$.
        \item Each question is answered at most once.
        \item The initial question is answered only once every non-initial question has been answered.
        \item For each P-answer $m_i$ there is an O-answer $m_j$ with the same label and $j < i$.
        \item The initial question is answered only if there is an O-answer to a root vertex with the same label as the initial question, or an initial assumption with the same label as the initial question.
    \end{enumerate}
    The datum of the initial question and of the initial assumptions is an atom; the datum of every other move is a node of the arena. The \emph{label} of a move is
its atom, or the atom labelling its node. \fillBox
\end{definition}

\begin{definition}[Strategy]
\label{def:strategy}
A \emph{strategy} for an arena $\arena{A}$ is a function $\sigma$ mapping
each play $\pi = (m_0,\ldots,m_k)$ whose last move is an O-move to
a (possibly empty) sequence of moves $\sigma(\pi) = (m_{k+1}, \ldots, m_n)$,
subject to the following conditions:
\begin{itemize}[align=left]
    \item \emph{(coherence)} the concatenation
          $\pi \cdot \sigma(\pi)$ is a play;
    \item \emph{(responsiveness)} if $\pi$ contains no unanswered P-questions that could be answered by Opponent in the next move, then $\sigma(\pi)$ is non-empty.
\end{itemize}
A strategy $\sigma$ is \emph{winning} for an atom $r$ from a set of
assumptions $Q$ if it satisfies the following further condition:
\begin{itemize}
    \item \emph{(winning)} every maximal play that opens with the initial
          question O?$r$ followed by the initial assumptions O!$q$, for
          $q \in Q$, and in which Proponent plays according to $\sigma$,
          contains a P-answer to the initial question.
\end{itemize}
We write $Q \wstrat{\arena{A}} r$ to express that there is a winning
strategy for $r$ from $Q$ in $\arena{A}$. \fillBox
\end{definition}

\begin{theorem}[Soundness for the game interpretation]
\label{thm:game-soundness}
    Let $\base{B}$ be a base, let $S$ be a finite set of atoms, and let $p$
    be an atom. If $S \deriv{\base{B}} p$, then
    $S \wstrat{\toarena{\base{B}}} p$.
\end{theorem}

\begin{proof}
By Definition~\ref{def:derivability-in-a-base}, the family of relations
$\deriv{\base{B}}$, indexed by bases $\base{B}$, is the least family closed
under (Ref) and (App). It therefore suffices to show that the family of
relations $\wstrat{\toarena{\base{B}}}$ is closed under the corresponding
conditions:
\[
\begin{array}{rl}
    \mbox{\rm (Ref)} & \mbox{if $p \in S$, then $S \wstrat{\toarena{\base{B}}} p$;} \\
    \mbox{\rm (App)} & \mbox{if $(Q_1 \brule q_1), \ldots, (Q_n \brule q_n)
        \brule r$ is a rule of $\base{B}$ and
        $S \wstrat{\toarena{\base{B} \cup Q_i}} q_i$,} \\
        & \mbox{for each $i = 1, \ldots, n$, then $S \wstrat{\toarena{\base{B}}} r$.}
\end{array}
\]
The theorem then follows by induction on the definition of
$S \deriv{\base{B}} p$.
 
\medskip
 
\noindent\emph{Case (Ref).} Suppose $p \in S$. By clauses~(i) and~(ii) of
Definition~\ref{def:play}, every play for $p$ from $S$ opens with the initial
question O?$p$ followed by the initial assumptions O!$s_1, \ldots,$ O!$s_k$,
where $s_1, \ldots, s_k$ enumerates $S$; since $p \in S$, the opening contains
the initial assumption O!$p$. Let $\sigma$ be the strategy that responds to
this opening with P!$p$. The move is legitimate: clause~(ix) is met by the
initial assumption O!$p$; clause~(x) permits answering the initial question,
again by that initial assumption; and clause~(viii) is satisfied vacuously,
the play containing no non-initial question. The resulting play
\[
\text{O?}p,\ \text{O!}s_1,\ \ldots,\ \text{O!}s_k,\ \text{P!}p
\]
is maximal --- the initial question has been answered, and by clause~(vii) no
question may be answered twice --- and it is the unique maximal play in which
Proponent follows $\sigma$. It contains a P-answer to the initial question,
so $\sigma$ is winning and $S \wstrat{\toarena{\base{B}}} p$.
 
\medskip
 
\noindent\emph{Case (App).} Suppose
$\rho = (Q_1 \brule q_1), \ldots, (Q_n \brule q_n) \brule r$ is a rule of
$\base{B}$ and, for each $i = 1, \ldots, n$, let $\sigma_i$ be a winning
strategy witnessing $S \wstrat{\toarena{\base{B} \cup Q_i}} q_i$. We
construct a winning strategy $\sigma$ witnessing
$S \wstrat{\toarena{\base{B}}} r$. We first record the shape of the arenas involved. The arena $\toarena{\{\rho\}}$ is a single
tree $\Psi$ whose root is labelled $r$, whose children are labelled
$q_1, \ldots, q_n$, and in which the subtrees below $q_i$ are exactly the
trees $\toarena{Q_i}$. Observe that
$\Psi$ is an element of $\toarena{\base{B}}$ and
$\toarena{\base{B} \cup Q_i} = \toarena{\base{B}} \cup \toarena{Q_i}$.
 
By clauses~(i) and~(ii) of Definition~\ref{def:play}, every play for $r$ from
$S$ opens O?$r$, O!$s_1$, \ldots, O!$s_k$, with $s_1, \ldots, s_k$
enumerating $S$. The strategy $\sigma$ extends this opening with the root
question P?$r$ at the root of $\Psi$, which clause~(iv) licenses: its node is
a root of the forest $\toarena{\base{B}}$, and its justifying question is the
initial question. Opponent now has two kinds of move available: by
clause~(v), Opponent may challenge a premiss, playing O?$q_i$ at a child of
the root of $\Psi$, justified by P?$r$; by clause~(vi), Opponent may instead
concede, answering P?$r$ with O!$r$. (If $n = 0$, conceding is Opponent's
only option.) These remain Opponent's only options throughout, alongside the
moves arising inside the sub-plays described below.
 
Suppose Opponent challenges with O?$q_i$. Then $\sigma$ continues by playing
according to $\sigma_i$, under the following identification: regard O?$q_i$
as the initial question, and the initial assumptions
O!$s_1, \ldots,$ O!$s_k$ of the ambient play as the initial assumptions, of a
simulated play in the arena $\toarena{\base{B} \cup Q_i}$. Moves of
$\sigma_i$ at nodes of $\toarena{\base{B}}$ carry over unchanged, since
$\toarena{\base{B}}$ is a sub-forest of both arenas. Moves of $\sigma_i$ at
nodes of $\toarena{Q_i}$ are transported to the corresponding nodes of $\Psi$
below $q_i$. The only clause of Definition~\ref{def:play} sensitive to this
relocation is the justification of questions at the roots of
$\toarena{Q_i}$: in the simulated play these are root questions, justified by
the initial question O?$q_i$ via clause~(iv), while in the ambient play the
corresponding nodes are children of $q_i$ in $\Psi$, so the same questions
are justified by O?$q_i$ via clause~(v). In either case the justifying
question is O?$q_i$, and every other justification is preserved. Clause~(x)
constrains only the answer to the ambient initial question, so it imposes
nothing on the sub-plays; clauses~(vi), (vii), and~(ix) hold in the ambient
play because they hold, move for move, in the simulated one.
 
Opponent may interleave the challenges O?$q_1, \ldots,$ O?$q_n$ and the
sub-plays they open. No interference results: justification in
Definition~\ref{def:play} relates \emph{occurrences} of moves, so each
non-initial move of the ambient play other than P?$r$ and O!$r$ descends, by
tracing justifying questions and the responses of $\sigma$, from a unique
challenge O?$q_i$, and $\sigma$ answers each O-move according to the
sub-strategy $\sigma_i$ of the challenge from which it descends. The
sub-plays therefore proceed independently, each a play of $\sigma_i$ under
the identification above. Since $\sigma_i$ is winning, every maximal
continuation of the sub-play opened by O?$q_i$ contains a P-answer to its
simulated initial question
% --- that is, the move P!$q_i$ ---
(i.e., the move P!$q_i$)
and, since
$\sigma_i$ respects clause~(viii) in the simulated play, every question of
that sub-play has been answered by the time P!$q_i$ is played.
 
Finally, suppose Opponent concedes O!$r$ (whether immediately after P?$r$ or
after any number of challenges). Once every non-initial question of the play
has been answered, $\sigma$ answers the initial question with P!$r$. The move
is legitimate: clause~(ix) is met by O!$r$; clause~(x) is met because O!$r$
is an O-answer at a root vertex whose label coincides with that of the
initial question; and clause~(viii) holds by the hypothesis that no
non-initial question stands unanswered.
 
It remains to check that $\sigma$ is winning. Let $\pi$ be a maximal play in
which Proponent follows $\sigma$; recall that plays are finite sequences
(Definition~\ref{def:play}). The question P?$r$ must be answered in $\pi$:
while it stands unanswered, clause~(vi) provides Opponent with the legal move
O!$r$ whenever it is Opponent's turn, and responsiveness
(Definition~\ref{def:strategy}) ensures that the play cannot halt on
Proponent's turn; hence a play in which P?$r$ is unanswered is not maximal.
By the same reasoning, every challenge O?$q_i$ that Opponent raises is
answered with P!$q_i$, since $\sigma_i$ is winning, and every question
internal to a sub-play is answered before the corresponding P!$q_i$, as
observed above. Thus every non-initial question of $\pi$ is answered, O!$r$
occurs in $\pi$, and $\sigma$ concludes with P!$r$. Every maximal play in
which Proponent follows $\sigma$ therefore contains a P-answer to the initial
question, so $\sigma$ is winning and $S \wstrat{\toarena{\base{B}}} r$.
\end{proof}

\end{document}
